\documentclass[11pt]{article}

\usepackage[T1]{fontenc}
\usepackage{lmodern}
\usepackage{microtype}
\usepackage[letterpaper,margin=1in]{geometry}
\usepackage{amsmath,amssymb,amsthm,mathtools}
\usepackage{enumitem}
\usepackage[hidelinks]{hyperref}

\hypersetup{
  pdftitle={A factorially weighted constant-term theorem on algebraic tori and the Gaussian Moments Conjecture in dimension two},
  pdfauthor={Christopher D. Long},
  pdfsubject={Factorially weighted constant terms, Mathieu--Zhao subspaces, and Gaussian moments},
  pdfkeywords={Gaussian moments, Mathieu--Zhao subspace, Laurent polynomial, constant term, Newton polytope, nullcone, algebraic torus, reduction modulo a prime}
}

\numberwithin{equation}{section}

\newtheorem{theorem}{Theorem}[section]
\newtheorem{proposition}[theorem]{Proposition}
\newtheorem{lemma}[theorem]{Lemma}
\newtheorem{corollary}[theorem]{Corollary}
\newtheorem{conjecture}[theorem]{Conjecture}
\theoremstyle{definition}
\newtheorem{definition}[theorem]{Definition}
\theoremstyle{remark}
\newtheorem{remark}[theorem]{Remark}

\newcommand{\C}{\mathbb C}
\newcommand{\R}{\mathbb R}
\newcommand{\Z}{\mathbb Z}
\newcommand{\N}{\mathbb N}
\newcommand{\T}{\mathbb T}
\newcommand{\E}{\mathbb E}
\newcommand{\A}{\mathcal A}
\newcommand{\supp}{\operatorname{supp}}
\newcommand{\conv}{\operatorname{conv}}
\newcommand{\CT}{\operatorname{CT}}
\newcommand{\Spec}{\operatorname{Spec}}
\newcommand{\GMC}{\operatorname{GMC}}
\newcommand{\SO}{\operatorname{SO}}
\newcommand{\Span}{\operatorname{span}}

\title{A Factorially Weighted Constant-Term Theorem on Algebraic Tori\\
\large and the Gaussian Moments Conjecture in Dimension Two}
\author{Christopher D. Long\\[2pt]
\href{mailto:galizur@gmail.com}{\texttt{galizur@gmail.com}}}
\date{July 22, 2026}

\begin{document}

\maketitle

\begin{abstract}
For an integer $r\ge 1$, let
\[
  \A_r=\C[u,z_1^{\pm1},\ldots,z_r^{\pm1}]
\]
and define the factorially weighted constant-term functional by
\[
  \Gamma_r(u^n z^\mu)=n!\,\mathbf 1_{\mu=0}
  \qquad(n\ge0,\ \mu\in\Z^r).
\]
For $F\in\A_r$, let $\Omega(F)\subset\Z^r$ be its torus-weight support.  We
prove that the following are equivalent: $\Gamma_r(F^m)=0$ for all
sufficiently large $m$; $\Gamma_r(F^m)=0$ for every $m\ge1$; and
$0\notin\conv\Omega(F)$.  Under these conditions,
$\Gamma_r(GF^m)=0$ for every fixed $G\in\A_r$ and all sufficiently large
$m$.  Hence $\ker\Gamma_r$ is a Mathieu--Zhao subspace for every torus rank.
More quantitatively, if $0\in\conv\Omega(F)$, then there is an integer
$k\ge1$ such that $\Gamma_r(F^{kp})\ne0$ for every sufficiently large prime
$p$.  The proof combines the multivariable constant-term theorem of
Duistermaat and van der Kallen, an exposed-face construction, and a
reduction-modulo-$p$ argument that isolates the lowest balanced coefficient
through factorial weights.  As a specialization, an explicit embedding of
the two-variable Gaussian functional into $\Gamma_1$ proves the Gaussian
Moments Conjecture in dimension two and the corresponding orthogonal
nullcone classification.
\end{abstract}

\medskip
\noindent\textbf{2020 Mathematics Subject Classification.}
Primary 13A50; Secondary 14L24, 14R15, 52B20, 60E05.

\smallskip
\noindent\textbf{Keywords.}
Factorially weighted constant term, algebraic torus, Gaussian moments,
Mathieu--Zhao subspace, Laurent polynomial, Newton polytope, nullcone,
reduction modulo a prime.

\section*{AI-assisted research disclaimer}
The results in this manuscript were generated by an interactive session
between ChatGPT 5.6 Sol and the author, Christopher D. Long.  The results
have not yet been human verified or formalized, but are believed to be
correct.  ChatGPT 5.6 Sol assisted with the drafting and editing of the
manuscript, and Claude Fable 5 supplied additional auditing and comments.

\section{Introduction}

For $r\ge1$, consider the Laurent polynomial algebra with one distinguished
nonnegative radial variable
\[
  \A_r=\C[u,z_1^{\pm1},\ldots,z_r^{\pm1}].
\]
Writing $z^\mu=z_1^{\mu_1}\cdots z_r^{\mu_r}$ for
$\mu=(\mu_1,\ldots,\mu_r)\in\Z^r$, define a linear functional
\[
  \Gamma_r:\A_r\longrightarrow\C,
  \qquad
  \Gamma_r(u^n z^\mu)=n!\,\mathbf 1_{\mu=0}.
\]
Equivalently, with normalized Haar measure on $\T^r$,
\begin{equation}\label{eq:gamma-integral-intro}
  \Gamma_r(F)
  =\int_0^\infty\int_{\T^r}F(u,z)e^{-u}\,dz\,du.
\end{equation}
Thus $\Gamma_r$ first extracts torus weight zero and then applies the Gamma
moment $u^n\mapsto n!$.

For
\[
  F=\sum_{\substack{n\ge0\\ \mu\in\Z^r}}c_{n,\mu}u^nz^\mu\in\A_r,
\]
define its torus-weight support by
\begin{equation}\label{eq:omega-intro}
  \Omega(F)
  =\{\mu\in\Z^r:c_{n,\mu}\ne0\text{ for some }n\ge0\}.
\end{equation}
We use the convention $\Omega(0)=\varnothing$ and
$\conv\varnothing=\varnothing$.

The main theorem classifies all polynomials whose sufficiently high power
moments vanish.

\begin{theorem}[Factorially weighted torus theorem]\label{thm:main}
Let $r\ge1$ and $F\in\A_r$.  The following conditions are equivalent.
\begin{enumerate}[label=\textup{(\roman*)}]
  \item $\Gamma_r(F^m)=0$ for all sufficiently large integers $m$;
  \item $\Gamma_r(F^m)=0$ for every integer $m\ge1$;
  \item $0\notin\conv\Omega(F)$.
\end{enumerate}
Moreover, whenever these conditions hold, for every $G\in\A_r$ one has
\[
  \Gamma_r(GF^m)=0\qquad(m\gg0).
\]
\end{theorem}

The quantitative form is stronger than the failure of eventual vanishing.

\begin{theorem}[Prime-index nonvanishing]\label{thm:prime-nonvanishing}
Let $r\ge1$ and let $F\in\A_r$ be nonzero.  If
$0\in\conv\Omega(F)$, then there exist integers $k\ge1$ and $p_0\ge1$ such
that, for every rational prime $p>p_0$,
\[
  \Gamma_r(F^{kp})\ne0.
\]
\end{theorem}

Recall that a linear subspace $M$ of a commutative $\C$-algebra $B$ is a
\emph{Mathieu--Zhao subspace} if, whenever $a^m\in M$ for every $m\ge1$,
one has $ba^m\in M$ for all sufficiently large $m$, for every $b\in B$;
see \cite{Zhao2012}.  The mixed-moment assertion in
Theorem~\ref{thm:main} immediately gives the following consequence.

\begin{corollary}\label{cor:mz-intro}
For every $r\ge1$, the kernel of $\Gamma_r$ is a Mathieu--Zhao subspace of
$\A_r$.
\end{corollary}

The proof has two independent components.  First, if the origin lies in the
convex hull of the torus weights, a linear-programming argument exposes the
lowest face capable of producing total torus weight zero.  The theorem of
Duistermaat and van der Kallen then supplies a nonzero balanced coefficient
on a power of the face polynomial.  Second, after normalizing by the
factorial attached to this lowest coefficient, reduction modulo a prime
shows that all higher radial terms vanish while the distinguished term
survives.  The exposed-face argument works in arbitrary torus rank; the
arithmetic isolation remains one-dimensional because there is only one
radial exponent.

The two-dimensional Gaussian theorem is obtained from the rank-one case.
If $X,Y$ are independent standard real Gaussian random variables, set
\begin{equation}\label{eq:circular-intro}
  Z=\frac{X+iY}{\sqrt2},
  \qquad
  W=\frac{X-iY}{\sqrt2}.
\end{equation}
The Gaussian contraction rule is
\[
  \E[Z^aW^b]=a!\,\mathbf 1_{a=b}.
\]
The injective algebra homomorphism
\begin{equation}\label{eq:iota-intro}
  \iota:\C[Z,W]\hookrightarrow\A_1,
  \qquad
  \iota(Z)=z,
  \qquad
  \iota(W)=uz^{-1},
\end{equation}
intertwines Gaussian expectation with $\Gamma_1$.  Consequently,
Theorem~\ref{thm:main} yields a complete classification of Gaussian
moment-null polynomials in two variables, proves $\GMC(2)$, and verifies the
corresponding nullcone conjecture of Derksen, van den Essen, and Zhao.

\section{The Gamma functional and torus weights}

For the remainder of the paper, fix $r\ge1$.  Every element of $\A_r$ has a
unique finite expansion
\[
  F=\sum_{\substack{n\ge0\\ \mu\in\Z^r}}c_{n,\mu}u^nz^\mu.
\]
The functional is therefore
\begin{equation}\label{eq:gamma-coefficients}
  \Gamma_r(F)=\sum_{n\ge0}n!\,[u^nz^0]F.
\end{equation}
The integral representation \eqref{eq:gamma-integral-intro} follows from
\[
  \int_{\T^r}z^\mu\,dz=\mathbf 1_{\mu=0},
  \qquad
  \int_0^\infty u^ne^{-u}\,du=n!.
\]

The torus weights give a $\Z^r$-grading
\begin{equation}\label{eq:torus-grading}
  \A_r=\bigoplus_{\mu\in\Z^r}(\A_r)_\mu,
  \qquad
  (\A_r)_\mu=\Span_\C\{u^nz^\mu:n\ge0\}.
\end{equation}
By definition, $\Gamma_r$ vanishes on every nonzero weight space.

\begin{lemma}[Strict separation and mixed moments]\label{lem:separation}
Suppose that $F\in\A_r$ satisfies
$0\notin\conv\Omega(F)$.  Then
\[
  \Gamma_r(F^m)=0\qquad(m\ge1),
\]
and, for every $G\in\A_r$,
\[
  \Gamma_r(GF^m)=0\qquad(m\gg0).
\]
\end{lemma}

\begin{proof}
The assertion is immediate for $F=0$, so assume $F\ne0$.  Since
$\conv\Omega(F)$ is compact and does not contain the origin, strict
separation gives a linear functional $\lambda:\R^r\to\R$ and a number
$\delta>0$ such that
\[
  \lambda(\mu)\ge\delta
  \qquad(\mu\in\Omega(F)).
\]
Every torus weight occurring in $F^m$ therefore has $\lambda$-value at least
$m\delta$ and is nonzero.  Hence $F^m$ has no weight-zero component and
$\Gamma_r(F^m)=0$.

If $G\ne0$, let
\[
  c=\min\{\lambda(\eta):\eta\in\Omega(G)\}.
\]
Every weight in $GF^m$ has $\lambda$-value at least $c+m\delta$, which is
positive for all sufficiently large $m$.  Thus $GF^m$ has no weight-zero
component and $\Gamma_r(GF^m)=0$.  The case $G=0$ is trivial.
\end{proof}

\section{The lowest balanced exposed face}

For a Laurent polynomial
$f\in\C[z_1^{\pm1},\ldots,z_r^{\pm1}]$, write $\CT(f)$ for its constant
term.  We use the full torus form of the theorem of Duistermaat and van der
Kallen.

\begin{theorem}[Duistermaat--van der Kallen]\label{thm:dvdk}
Let $f\in\C[z_1^{\pm1},\ldots,z_r^{\pm1}]$.  If
\[
  \CT(f^m)=0\qquad\text{for every }m\ge1,
\]
then
\[
  0\notin\conv\supp(f).
\]
Equivalently, if $0\in\conv\supp(f)$, then $\CT(f^k)\ne0$ for at least one
integer $k\ge1$.
\end{theorem}

\begin{proof}[Reference]
This is \cite[Theorem~1]{DuistermaatVanderKallen1998}.
\end{proof}

We next record the supporting-face statement used to expose the least radial
degree compatible with total torus weight zero.

\begin{lemma}[Lowest balanced face]\label{lem:lowest-face}
Let $S$ be a finite set, let $w:S\to\R^r$, and let $n:S\to\R$.  Assume that
$0\in\conv w(S)$.  Define
\begin{equation}\label{eq:rho-primal}
  \rho=
  \min\left\{
    \sum_{s\in S}t_s n(s):
    t_s\ge0,\ \sum_{s\in S}t_s=1,\
    \sum_{s\in S}t_sw(s)=0
  \right\}.
\end{equation}
Then there exists $\nu\in\R^r$ such that
\begin{equation}\label{eq:supporting-inequality}
  n(s)+\langle\nu,w(s)\rangle\ge\rho
  \qquad(s\in S).
\end{equation}
If
\[
  S_0=\{s\in S:n(s)+\langle\nu,w(s)\rangle=\rho\},
\]
then $0\in\conv w(S_0)$.
\end{lemma}

\begin{proof}
The optimization problem \eqref{eq:rho-primal} is a feasible bounded linear
program.  Its dual is
\[
  \max\left\{\rho':
  \rho'\le n(s)+\langle\nu,w(s)\rangle
  \text{ for every }s\in S\right\},
\]
with variables $(\rho',\nu)\in\R\times\R^r$.  Strong linear-programming
duality gives an optimal pair $(\rho,\nu)$; see, for example,
\cite[Chapter~7]{Schrijver1986}.  If $(t_s)$ is an optimal solution of
\eqref{eq:rho-primal}, then
\[
  0=\sum_{s\in S}t_s
  \bigl(n(s)+\langle\nu,w(s)\rangle-\rho\bigr).
\]
Every summand is nonnegative, so $t_s>0$ only for $s\in S_0$.  Since the
$t_s$ sum to one and their $w$-weighted average is zero, it follows that
$0\in\conv w(S_0)$.
\end{proof}

\begin{proposition}[Balanced face extraction]\label{prop:face-extraction}
Let $F\in\A_r$ be nonzero and assume $0\in\conv\Omega(F)$.  Then there
exist
\begin{itemize}
  \item an integer $k\ge1$,
  \item an integer $d\ge0$,
  \item a vector $\nu\in\R^r$, and
  \item a nonzero complex number $C$,
\end{itemize}
such that, for $R=F^k$,
\begin{enumerate}[label=\textup{(\alph*)}]
  \item every monomial $u^Nz^\mu$ occurring in $R$ satisfies
  \begin{equation}\label{eq:ell-lower-bound}
    N+\langle\nu,\mu\rangle\ge d;
  \end{equation}
  \item
  \begin{equation}\label{eq:distinguished-coefficient}
    [u^dz^0]R=C\ne0.
  \end{equation}
\end{enumerate}
\end{proposition}

\begin{proof}
Write
\[
  F=\sum_{(n,\mu)\in S}c_{n,\mu}u^nz^\mu,
  \qquad
  S=\supp(F)\subset\Z_{\ge0}\times\Z^r.
\]
Apply Lemma~\ref{lem:lowest-face} with
\[
  w(n,\mu)=\mu,
  \qquad
  n(n,\mu)=n.
\]
Thus there are $\rho\in\R$ and $\nu\in\R^r$ such that
\begin{equation}\label{eq:ell-def}
  \ell(n,\mu):=n+\langle\nu,\mu\rangle\ge\rho
  \qquad((n,\mu)\in S),
\end{equation}
and the equality set
\[
  S_0=\{(n,\mu)\in S:\ell(n,\mu)=\rho\}
\]
satisfies $0\in\conv\{\mu:(n,\mu)\in S_0\}$.

Projection from $S_0$ to $\Z^r$ is injective.  Indeed, if $(n,\mu)$ and
$(n',\mu)$ both lie on the face, then
\[
  n+\langle\nu,\mu\rangle
  =n'+\langle\nu,\mu\rangle=\rho,
\]
so $n=n'$.  Consequently, for the face polynomial
\[
  H(u,z)=\sum_{(n,\mu)\in S_0}c_{n,\mu}u^nz^\mu
\]
and the Laurent polynomial
\begin{equation}\label{eq:face-laurent}
  h(z)=\sum_{(n,\mu)\in S_0}c_{n,\mu}z^\mu,
\end{equation}
there is no cancellation when equal Laurent monomials are collected, and
\[
  0\in\conv\supp(h).
\]
By the contrapositive of Theorem~\ref{thm:dvdk}, there is an integer $k\ge1$
such that
\begin{equation}\label{eq:C-constant-term}
  C:=\CT(h^k)\ne0.
\end{equation}

Every contribution to the constant term in \eqref{eq:C-constant-term}
comes from a product of $k$ face monomials whose total torus weight is zero.
Because every factor lies on $\ell=\rho$, additivity gives
\[
  \ell(N,0)=N=k\rho.
\]
Thus every balanced product of $k$ face monomials has the same radial degree
\[
  d:=k\rho\in\Z_{\ge0},
\]
and therefore
\begin{equation}\label{eq:H-coefficient}
  [u^dz^0]H^k=\CT(h^k)=C\ne0.
\end{equation}

Set $R=F^k$.  Every monomial of $R$ has $\ell$-degree at least $k\rho=d$.
A product in the expansion of $F^k$ that uses at least one monomial outside
$S_0$ has $\ell$-degree strictly greater than $d$ and hence cannot
contribute to $u^dz^0$.  Equation \eqref{eq:H-coefficient} therefore gives
$[u^dz^0]R=C$, and \eqref{eq:ell-lower-bound} follows from additivity.
\end{proof}

\section{Factorial isolation by reduction modulo a prime}

We first record a spreading-out fact in the form required below.

\begin{lemma}[Good reductions]\label{lem:good-reductions}
Let $B$ be a finitely generated integral domain over $\Z$ of characteristic
zero.  Then there is an integer $N\ge1$ such that, for every prime
$p\nmid N$, there exist a field $K_p$ of characteristic $p$ and a unital ring
homomorphism
\[
  \varphi_p:B\longrightarrow K_p.
\]
In particular, every unit of $B$ has nonzero image in $K_p$.
\end{lemma}

\begin{proof}
The morphism
\[
  \pi:\Spec(B)\longrightarrow\Spec(\Z)
\]
is dominant because the zero ideal of the domain $B$ contracts to the zero
ideal of $\Z$.  Since $\Z$ is Noetherian, the finite-type $\Z$-algebra $B$
is finitely presented.  Chevalley's theorem shows that the image of $\pi$ is
constructible; see \cite[Tag 00FE]{StacksProject}.

This image contains the generic point of the irreducible Noetherian space
$\Spec(\Z)$.  A constructible subset of an irreducible Noetherian space that
contains the generic point contains a nonempty open subset.  It follows that
the image of $\pi$ contains $D(N)$ for some integer $N\ge1$.

For every prime $p\nmid N$, choose a prime ideal $\mathfrak q\subset B$ with
$\mathfrak q\cap\Z=(p)$ and set
\[
  K_p=\operatorname{Frac}(B/\mathfrak q).
\]
Then $K_p$ has characteristic $p$, and the composite
$B\to B/\mathfrak q\hookrightarrow K_p$ is the required homomorphism.  A
unit cannot have zero image under a unital homomorphism to a field.
\end{proof}

The next lemma is the arithmetic core of the argument.

\begin{lemma}[Factorial isolation]\label{lem:factorial-isolation}
Let
\[
  R=\sum_{\substack{n\ge0\\ \mu\in\Z^r}}
  r_{n,\mu}u^nz^\mu\in\A_r.
\]
Suppose that there are $\nu\in\R^r$ and an integer $d\ge0$ such that
\begin{equation}\label{eq:R-lower-degree}
  n+\langle\nu,\mu\rangle\ge d
  \qquad\text{whenever }r_{n,\mu}\ne0,
\end{equation}
and suppose
\begin{equation}\label{eq:C-rd0}
  C:=r_{d,0}\ne0.
\end{equation}
Then
\[
  \Gamma_r(R^p)\ne0
\]
for every sufficiently large prime $p$.
\end{lemma}

\begin{proof}
Let
\begin{equation}\label{eq:coefficient-ring}
  B=\Z[\{r_{n,\mu}\},C^{-1}]\subset\C,
\end{equation}
where only the finitely many nonzero coefficients of $R$ are included.  This
is a finitely generated integral domain of characteristic zero.  Choose $N$
as in Lemma~\ref{lem:good-reductions}.

Fix a prime $p\nmid N$.  Every torus-weight-zero monomial $u^jz^0$ occurring
in $R^p$ satisfies
\[
  j=j+\langle\nu,0\rangle\ge pd,
\]
because the degree $(n,\mu)\mapsto n+\langle\nu,\mu\rangle$ is additive and
all monomials of $R$ have degree at least $d$.  Define
\begin{equation}\label{eq:Gp-def}
  G_p=
  \sum_{j\ge pd}[u^jz^0]R^p\,\frac{j!}{(pd)!}\in B.
\end{equation}
By \eqref{eq:gamma-coefficients},
\begin{equation}\label{eq:moment-Gp}
  \Gamma_r(R^p)=(pd)!\,G_p.
\end{equation}

Let $\varphi_p:B\to K_p$ be supplied by
Lemma~\ref{lem:good-reductions}, and place a bar over reductions to $K_p$.
In characteristic $p$, the Frobenius identity gives
\begin{equation}\label{eq:frobenius}
  \overline{R^p}=\bar R^{\,p}
  =\sum_{n,\mu}\bar r_{n,\mu}^{\,p}u^{pn}z^{p\mu}.
\end{equation}
Only terms with $\mu=0$ can have torus weight zero after reduction, and
therefore
\begin{equation}\label{eq:Gp-reduction}
  \overline{G_p}
  =\sum_{n\ge d}\bar r_{n,0}^{\,p}
    \frac{(pn)!}{(pd)!}
  \quad\text{in }K_p.
\end{equation}
For $n=d$, the factorial ratio is $1$, and the corresponding term is
$\bar C^{\,p}$.  If $n>d$, then
\[
  \frac{(pn)!}{(pd)!}
  =(pd+1)(pd+2)\cdots(pn)
\]
is divisible by $p$, since the displayed product contains the factor
$p(d+1)$.  Hence all terms with $n>d$ vanish in $K_p$, while
\begin{equation}\label{eq:Gp-isolated}
  \overline{G_p}=\bar C^{\,p}\ne0.
\end{equation}
The last inequality holds because $C$ is a unit in $B$.  Thus $G_p\ne0$,
and \eqref{eq:moment-Gp} implies $\Gamma_r(R^p)\ne0$.  This holds for every
prime $p\nmid N$, hence for every sufficiently large prime.
\end{proof}

\begin{remark}
The normalization by $(pd)!$ is essential: it converts the higher radial
terms into factorial ratios divisible by $p$, while the distinguished lowest
balanced coefficient retains ratio $1$.
\end{remark}

\section{Proof of the main theorem}

\begin{proof}[Proof of Theorem~\ref{thm:prime-nonvanishing}]
Apply Proposition~\ref{prop:face-extraction} to obtain
$R=F^k$, $d$, $\nu$, and $C\ne0$ satisfying
\eqref{eq:ell-lower-bound} and \eqref{eq:distinguished-coefficient}.
Lemma~\ref{lem:factorial-isolation} gives
\[
  \Gamma_r(F^{kp})=\Gamma_r(R^p)\ne0
\]
for every sufficiently large prime $p$.
\end{proof}

\begin{proof}[Proof of Theorem~\ref{thm:main}]
The implication \textup{(ii)}$\Rightarrow$\textup{(i)} is immediate.
Lemma~\ref{lem:separation} proves
\textup{(iii)}$\Rightarrow$\textup{(ii)} and also proves the mixed-moment
conclusion.

It remains to prove \textup{(i)}$\Rightarrow$\textup{(iii)}.  The case
$F=0$ is immediate.  Suppose $F\ne0$ and
$\Gamma_r(F^m)=0$ for all sufficiently large $m$.  If
$0\in\conv\Omega(F)$, Theorem~\ref{thm:prime-nonvanishing} gives an integer
$k\ge1$ such that
\[
  \Gamma_r(F^{kp})\ne0
\]
for every sufficiently large prime $p$, contradicting eventual vanishing.
Therefore $0\notin\conv\Omega(F)$.
\end{proof}

\begin{proof}[Proof of Corollary~\ref{cor:mz-intro}]
Suppose $F^m\in\ker\Gamma_r$ for every $m\ge1$.  By
Theorem~\ref{thm:main}, $0\notin\conv\Omega(F)$.  The mixed-moment part of
the theorem gives $\Gamma_r(GF^m)=0$ for all sufficiently large $m$ and
every $G\in\A_r$.  This is precisely the Mathieu--Zhao property.
\end{proof}

Theorem~\ref{thm:main} also identifies the radical of the kernel in support
terms.

\begin{corollary}[Radical classification]\label{cor:radical}
For every $r\ge1$,
\[
  \{F\in\A_r:\Gamma_r(F^m)=0\text{ for }m\gg0\}
  =\{F\in\A_r:0\notin\conv\Omega(F)\}.
\]
If $F\ne0$ belongs to this set, there is a linear functional
$\lambda:\R^r\to\R$ that is strictly positive on every weight of $F$.
\end{corollary}

\begin{proof}
The equality is Theorem~\ref{thm:main}.  The final assertion is the strict
separation theorem applied to the compact polytope $\conv\Omega(F)$ and the
origin.
\end{proof}

\section{The two-dimensional Gaussian specialization}

Let $X,Y$ be independent standard real Gaussian random variables and define
$Z,W$ by \eqref{eq:circular-intro}.  Since the change of variables is
invertible over $\C$, we identify $\C[X,Y]$ with $\C[Z,W]$.

\begin{lemma}[Gaussian contraction rule]\label{lem:gaussian-contraction}
For all $a,b\in\Z_{\ge0}$,
\begin{equation}\label{eq:gaussian-contraction}
  \E[Z^aW^b]
  =\begin{cases}
     a!,&a=b,\\
     0,&a\ne b.
   \end{cases}
\end{equation}
\end{lemma}

\begin{proof}
Write $X=r\cos\theta$ and $Y=r\sin\theta$.  Under standard Gaussian
measure, $\theta$ is uniform on $[0,2\pi)$ and independent of
$U=r^2/2$, while $U$ has density $e^{-u}\,du$ on $[0,\infty)$.  Moreover,
\[
  Z=\sqrt U\,e^{i\theta},
  \qquad
  W=\sqrt U\,e^{-i\theta}.
\]
The angular integral vanishes unless $a=b$.  In that case,
\[
  \E[Z^aW^a]=\E[U^a]
  =\int_0^\infty u^ae^{-u}\,du=a!.
\]
\end{proof}

\begin{proposition}[Intertwining embedding]\label{prop:intertwining}
The assignment
\[
  \iota(Z)=z,
  \qquad
  \iota(W)=uz^{-1}
\]
defines an injective algebra homomorphism
$\iota:\C[Z,W]\hookrightarrow\A_1$ satisfying
\begin{equation}\label{eq:intertwining}
  \Gamma_1(\iota(P))=\E[P(Z,W)]
  \qquad(P\in\C[Z,W]).
\end{equation}
Moreover, the torus-weight support of $\iota(P)$ is exactly the angular
weight support
\[
  \Omega_{\mathrm{ang}}(P)
  =\{a-b:[Z^aW^b]P\ne0\}.
\]
\end{proposition}

\begin{proof}
For a monomial,
\[
  \iota(Z^aW^b)=u^bz^{a-b}.
\]
The pairs $(b,a-b)$ determine $(a,b)$ uniquely, so $\iota$ is injective.
Furthermore,
\[
  \Gamma_1(u^bz^{a-b})
  =b!\,\mathbf 1_{a=b}
  =a!\,\mathbf 1_{a=b}
  =\E[Z^aW^b]
\]
by Lemma~\ref{lem:gaussian-contraction}.  Linearity proves
\eqref{eq:intertwining}, and the support statement is immediate.
\end{proof}

\begin{theorem}[Two-variable Gaussian classification]\label{thm:gaussian-main}
Let $P\in\C[X,Y]$, viewed in circular coordinates as an element of
$\C[Z,W]$.  The following are equivalent.
\begin{enumerate}[label=\textup{(\roman*)}]
  \item $\E[P(X,Y)^m]=0$ for all sufficiently large $m$;
  \item $\E[P(X,Y)^m]=0$ for every $m\ge1$;
  \item every angular weight $a-b$ occurring in $P$ is strictly positive,
        or every such weight is strictly negative.
\end{enumerate}
Moreover, whenever these conditions hold, for every $Q\in\C[X,Y]$,
\[
  \E[Q(X,Y)P(X,Y)^m]=0\qquad(m\gg0).
\]
If the angular support of a nonzero $P$ contains zero in its convex hull,
then there is a $k\ge1$ such that
\[
  \E[P(X,Y)^{kp}]\ne0
\]
for every sufficiently large prime $p$.
\end{theorem}

\begin{proof}
Apply Theorems~\ref{thm:main} and \ref{thm:prime-nonvanishing} to
$\iota(P)$ and use Proposition~\ref{prop:intertwining}.  Since the angular
weights lie in the one-dimensional lattice $\Z$, exclusion of zero from
their convex hull is equivalent to all occurring weights having one strict
sign.
\end{proof}

\begin{corollary}[Gaussian Moments Conjecture in dimension two]
The Gaussian Moments Conjecture $\GMC(2)$ is true.
\end{corollary}

\begin{proof}
If $\E[P^m]=0$ for every $m\ge1$, the mixed-moment conclusion in
Theorem~\ref{thm:gaussian-main} gives
\[
  \E[QP^m]=0\qquad(m\gg0)
\]
for every polynomial $Q$.
\end{proof}

We finish by recording the orthogonal nullcone interpretation.  In circular
coordinates, consider the one-parameter subgroup
\begin{equation}\label{eq:lambda-ZW}
  (Z,W)\longmapsto(tZ,t^{-1}W).
\end{equation}
In the original coordinates it is
\begin{equation}\label{eq:lambda-XY}
  \lambda(t)=\frac12
  \begin{pmatrix}
    t+t^{-1} & i(t-t^{-1})\\
    -i(t-t^{-1}) & t+t^{-1}
  \end{pmatrix}.
\end{equation}
It preserves $X^2+Y^2=2ZW$ and has determinant one, so
$\lambda(t)\in\SO_2(\C)$.

\begin{corollary}[Orthogonal nullcone in dimension two]
Let $P\in\C[X,Y]$ satisfy
$\E[P(X,Y)^m]=0$ for every $m\ge1$.  Then there is a one-parameter subgroup
$\lambda:\C^\times\to\SO_2(\C)$ such that
\[
  P(\lambda(t)(X,Y))\in t\C[t][X,Y].
\]
Thus the nullcone conjecture of Derksen, van den Essen, and Zhao holds in
dimension two.
\end{corollary}

\begin{proof}
By Theorem~\ref{thm:gaussian-main}, the angular weights of $P$ are all
positive or all negative.  In the positive case, substitution
\eqref{eq:lambda-ZW} multiplies $Z^aW^b$ by $t^{a-b}$, and every exponent is
at least one.  In the negative case, use the inverse one-parameter subgroup.
The zero polynomial is immediate.
\end{proof}

\begin{thebibliography}{99}

\bibitem{DerksenEssenZhao2017}
H.~Derksen, A.~van den Essen, and W.~Zhao,
\emph{The Gaussian Moments Conjecture and the Jacobian Conjecture},
Israel J. Math. \textbf{219} (2017), no.~2, 917--928.
\href{https://doi.org/10.1007/s11856-017-1502-2}{doi:10.1007/s11856-017-1502-2}.

\bibitem{DuistermaatVanderKallen1998}
J.~J.~Duistermaat and W.~van der Kallen,
\emph{Constant terms in powers of a Laurent polynomial},
Indag. Math. (N.S.) \textbf{9} (1998), no.~2, 221--231.
\href{https://doi.org/10.1016/S0019-3577(98)80020-7}{doi:10.1016/S0019-3577(98)80020-7}.
MR1691479.

\bibitem{Schrijver1986}
A.~Schrijver,
\emph{Theory of Linear and Integer Programming},
Wiley-Interscience Series in Discrete Mathematics,
John Wiley \& Sons, Chichester, 1986.

\bibitem{StacksProject}
The Stacks Project Authors,
\emph{The Stacks Project},
Tag 00FE (Chevalley's theorem),
\url{https://stacks.math.columbia.edu/tag/00FE}.

\bibitem{Zhao2012}
W.~Zhao,
\emph{Mathieu subspaces of associative algebras},
J. Algebra \textbf{350} (2012), 245--272.
\href{https://doi.org/10.1016/j.jalgebra.2011.09.036}{doi:10.1016/j.jalgebra.2011.09.036}.
MR2859886.

\end{thebibliography}

\end{document}
